\section{Step 70 Require end-of-block before building the data tables}\label{png-huff-s70}
Fill \texttt{huff\_end\_code}. After exactly the declared combined entries have been expanded, require a nonzero literal/length length for symbol 256. Then follow the supplied calls that build the two data tables and start ordinary symbol decoding.

The dynamic-literals fixture assigns symbol zero length one, 256 and 257 length two, and its single distance symbol length one. Its data emits 35 zeros and then end-of-block. The complete inflater now returns zero and 35 for this valid stream and also retains the stored and fixed paths. The whole program still reports development status because reconstruction and packing are the next unfinished stages.

\paragraph{Observe} Use \texttt{s70-1}; stop before \texttt{huff\_end\_code\_ready}. Inspect the state described above before continuing.

\Needspace{7\baselineskip}
\paragraph{Check} Why require a code for symbol 256?
\begin{itemize}
  \item[A.] It supplies the image width.
  \item[B.] Every compressed block needs an end-of-block symbol.
  \item[C.] It always emits 256 bytes.
\end{itemize}
